% !TeX root = ../main.tex
% ---------------------------------------------------------------------------
% Seismic Design of Underground Structures
% Scope extension: not derived from the RSE3010 weekly material.
% Key sources: Wang (1993); Penzien (2000); Hashash et al. (2001) state of
% the art; FHWA Technical Manual for Design and Construction of Road Tunnels;
% AS 1170.4 for the Australian hazard context.
% ---------------------------------------------------------------------------
\chapter{Seismic Design of Underground Structures}
\label{ch:seismic-design}


\section{Seismic performance of underground structures}
\label{sec:seismic-performance-of-underground-structures}

\subsection{Observed seismic performance of tunnels}
\label{sec:seismic-performance}
\scopenote{Empirical damage databases, the general robustness of deep bored
tunnels, and the conditions under which that robustness fails --- shallow
cover, portals, fault crossings and liquefiable ground.}

\subsection{Ground shaking versus ground failure}
\label{sec:shaking-vs-failure}
\scopenote{Transient deformation imposed by passing waves, distinguished from
permanent displacement due to fault rupture, liquefaction or landsliding. The
two require entirely different design responses.}

\section{Hazard and deformation demand}
\label{sec:hazard-and-deformation-demand}

\subsection{Seismic hazard and design ground motion}
\label{sec:seismic-hazard}
\scopenote{Probabilistic hazard assessment, return periods for operating and
maximum design earthquakes, and the Australian intraplate hazard context.}

\subsection{Deformation modes of a linear structure}
\label{sec:deformation-modes}
\scopenote{Axial extension and compression, longitudinal bending, and
transverse ovaling or racking; the free-field strain from wave propagation.}

\section{Simplified analysis methods}
\label{sec:simplified-analysis-methods}

\subsection{Free-field deformation methods}
\label{sec:free-field}
\scopenote{Newmark's simplified wave-passage approach for axial and curvature
strain, angle of incidence assumptions, and the conservatism involved.}

\subsection{Ovaling and racking analysis}
\label{sec:ovaling-racking}
\scopenote{Wang's and Penzien's closed-form solutions for circular linings and
rectangular boxes; the compressibility and flexibility ratios that control
whether the lining follows or resists the ground.}

\subsection{Soil-structure interaction effects}
\label{sec:seismic-ssi}
\scopenote{Full-slip and no-slip interface assumptions, their effect on thrust
and moment, and when interaction analysis is required rather than free-field
imposition.}

\section{Dynamic numerical analysis}
\label{sec:dynamic-analysis}
\scopenote{Input motion selection and deconvolution, boundary conditions
(free-field, absorbing), damping, and nonlinear site response.}

\section{Permanent ground displacement}
\label{sec:permanent-ground-displacement}

\subsection{Fault crossing design}
\label{sec:fault-crossing}
\scopenote{Estimating displacement demand, oversized sections, articulated and
sacrificial lining segments, ductile detailing, and repairability as a design
objective.}

\subsection{Liquefaction and its consequences}
\label{sec:liquefaction}
\scopenote{Uplift and flotation of shallow tunnels, loss of lateral support,
and mitigation by ground improvement or increased weight.}

\section{Structures and performance criteria}
\label{sec:structures-and-performance-criteria}

\subsection{Portals, shafts and stations}
\label{sec:seismic-portals}
\scopenote{Why the transition zones and box structures rather than the bored
tunnel usually govern seismic design; slope stability at portals under seismic
loading.}

\subsection{Performance-based design criteria}
\label{sec:performance-based}
\scopenote{Performance levels from immediate operation to life safety, and
matching them to hazard levels and asset criticality.}
